\section{Background and Problem Definition}

The paper uses CVE IDs as the record-alignment key \cite{CVE2026} and studies field values drawn from NVD \cite{NVDAPI2026} and GHSA \cite{GHSA2026}. NVD enrichment adds analysis beyond the base CVE record, including CVSS metrics, CWE mappings, reference tagging, and vulnerable-product applicability information \cite{NVDProcess2026,NVDStatus2026}. GHSA combines GitHub-reviewed advisories, unreviewed community advisories, and NVD-imported records, and it exposes advisory data through an OSV-compatible JSON representation for supported ecosystems \cite{GHSAAbout2026,GHSARepo2026}. Severity follows CVSS-style score/vector representations \cite{CVSS31}, weakness categories are represented through CWE identifiers \cite{CWE2026}, and affected-version ranges are close to the OSV/GHSA advisory schema family \cite{OSVSchema2026}. These sources are not interchangeable: they differ in curation workflow, schema choices, publication timing, and the amount of advisory context exposed through references.

Affected-version comparison is especially schema-sensitive. OSV separates package ecosystem/name identity from affected version lists and range events, including introduced/fixed event boundaries, and GHSA maintains repository-specific extensions for version-range information that is not always expressible in plain OSV fields \cite{OSVSchema2026,GHSARepo2026}. The affected\_versions field is therefore treated as a structured vulnerability-record field whose schema includes package and range information; the current prototype uses token-support and structural diagnostics rather than semantic version-range adjudication.

Published/date discrepancies are also process-sensitive. NVD enrichment, GHSA review, and NVD-imported advisory flows can expose different timestamps for the same CVE because records move through different publication and update paths \cite{NVDProcess2026,NVDOperations2026,Segal2026GHSA}. The temporal\_discrepancy category is intended to keep those timing differences separate from substantive factual conflicts.

We treat an aligned NVD-GHSA pair as input and a field as the unit of analysis. For an aligned pair \texttt{(r\_NVD, r\_GHSA)} and field \texttt{f}, the first task is to assign a discrepancy type to \texttt{(r\_NVD.f, r\_GHSA.f)}. The second task is activated only for factual conflicts: given field values and fetched advisory evidence, record whether the available evidence cues match NVD, GHSA, both sources, neither source, or whether the case should abstain for audit.

The method uses five working labels: equivalent, representation\_discrepancy, incomplete, temporal\_discrepancy, and factual\_conflict. Equivalent means the normalized values carry the same field meaning. Representation\_discrepancy means the raw strings or schemas differ but the intended fact is treated as compatible by the deterministic rule. Incomplete means one side contains field content that the other side lacks. Temporal\_discrepancy captures publication/date differences that are better explained as timing differences than as contradictory facts. Factual\_conflict is reserved for residual cases where the implemented rules do not find compatibility, incompleteness, or temporal explanation.

This distinction matters because downstream actions differ by type, and prior work shows that vulnerability-data quality, source consistency, and scanner behavior affect operational use \cite{NVDUsers2024,Churakova2026VEX,Mandl2026VulnDetection}. Equivalent and representation discrepancies can usually be normalized away. Incompleteness suggests augmentation rather than source selection. Temporal discrepancies require timeline interpretation. Factual conflicts require evidence and may still end in abstention if the evidence is weak or contradictory.

The current labels are useful for organizing the analysis, but they are not gold truth. The method therefore avoids treating all field mismatches as factual conflicts or collapsing record alignment into the discrepancy task. The working problem is two-stage: first normalize and classify the field disagreement; then, if the field is a factual conflict, use evidence from references and related advisory material to decide whether one side is better supported. The current implementation records this problem definition and prototype; it does not claim that the adjudication stage is final.
